\documentclass[../20232006p.tex]{subfiles}
\begin{document}

Phụ lục này trình bày đặc tả chi tiết các Use Case của hệ thống quản lý quán cơm bình dân. Mỗi Use Case được mô tả theo chuẩn gồm: Mã định danh, Tên, Tác nhân, Mô tả, Điều kiện tiên quyết, Luồng chính và Luồng thay thế.

% ===================== KHÁCH HÀNG =====================
\section{Đặc tả Use Case phân hệ Khách hàng}

\subsection{UC-C01: Đặt hàng trực tuyến}
\begin{table}[H]
    \centering
    \begin{tabular}{|p{3.5cm}|p{10cm}|}
        \hline
        \textbf{Mã UC}                & UC-C01                                                                          \\
        \hline
        \textbf{Tên}                  & Đặt hàng trực tuyến                                                             \\
        \hline
        \textbf{Tác nhân}             & Khách hàng                                                                      \\
        \hline
        \textbf{Mô tả}                & Khách hàng đặt món ăn qua website, chọn thời gian lấy hàng và thanh toán online \\
        \hline
        \textbf{Điều kiện tiên quyết} & Khách hàng đã đăng nhập vào hệ thống                                            \\
        \hline
        \textbf{Luồng chính}          &
        1. Khách hàng xem thực đơn hôm nay \newline
        2. Khách hàng chọn món và thêm vào giỏ hàng \newline
        3. Khách hàng xem giỏ hàng và điều chỉnh số lượng \newline
        4. Khách hàng chọn địa chỉ giao hàng hoặc thời gian đến lấy \newline
        5. Khách hàng thanh toán qua VietQR \newline
        6. Hệ thống tạo đơn hàng và gửi thông báo xác nhận                                                              \\
        \hline
        \textbf{Luồng thay thế}       &
        3a. Giỏ hàng trống: Hệ thống hiển thị thông báo yêu cầu thêm món \newline
        5a. Thanh toán thất bại: Hệ thống hiển thị thông báo lỗi                                                        \\
        \hline
    \end{tabular}
\end{table}

\subsection{UC-C02: Xem lịch sử mua hàng}
\begin{table}[H]
    \centering
    \begin{tabular}{|p{3.5cm}|p{10cm}|}
        \hline
        \textbf{Mã UC}                & UC-C02                                       \\
        \hline
        \textbf{Tên}                  & Xem lịch sử mua hàng                         \\
        \hline
        \textbf{Tác nhân}             & Khách hàng                                   \\
        \hline
        \textbf{Mô tả}                & Khách hàng xem danh sách các đơn hàng đã đặt \\
        \hline
        \textbf{Điều kiện tiên quyết} & Khách hàng đã đăng nhập                      \\
        \hline
        \textbf{Luồng chính}          &
        1. Khách hàng chọn menu Lịch sử đơn hàng \newline
        2. Hệ thống hiển thị danh sách đơn hàng theo thời gian \newline
        3. Khách hàng có thể lọc theo trạng thái hoặc khoảng ngày \newline
        4. Khách hàng chọn xem chi tiết một đơn hàng                                 \\
        \hline
    \end{tabular}
\end{table}

\subsection{UC-C03: Gửi Feedback}
\begin{table}[H]
    \centering
    \begin{tabular}{|p{3.5cm}|p{10cm}|}
        \hline
        \textbf{Mã UC}                & UC-C03                                                     \\
        \hline
        \textbf{Tên}                  & Gửi Feedback                                               \\
        \hline
        \textbf{Tác nhân}             & Khách hàng                                                 \\
        \hline
        \textbf{Mô tả}                & Khách hàng đánh giá chất lượng đơn hàng sau khi hoàn thành \\
        \hline
        \textbf{Điều kiện tiên quyết} & Đơn hàng đã hoàn thành                                     \\
        \hline
        \textbf{Luồng chính}          &
        1. Khách hàng vào chi tiết đơn hàng đã hoàn thành \newline
        2. Khách hàng chọn số sao đánh giá (1-5) \newline
        3. Khách hàng nhập nhận xét (tùy chọn) \newline
        4. Hệ thống lưu feedback và hiển thị thông báo cảm ơn                                      \\
        \hline
    \end{tabular}
\end{table}

% ===================== NHÂN VIÊN =====================
\section{Đặc tả Use Case phân hệ Nhân viên bán hàng}

\subsection{UC-S01: Thanh toán đơn hàng}
\begin{table}[H]
    \centering
    \begin{tabular}{|p{3.5cm}|p{10cm}|}
        \hline
        \textbf{Mã UC}                & UC-S01                                                  \\
        \hline
        \textbf{Tên}                  & Thanh toán đơn hàng                                     \\
        \hline
        \textbf{Tác nhân}             & Nhân viên thu ngân                                      \\
        \hline
        \textbf{Mô tả}                & Nhân viên tạo và thanh toán đơn hàng tại quầy cho khách \\
        \hline
        \textbf{Điều kiện tiên quyết} & Nhân viên đã đăng nhập và mở ca làm việc                \\
        \hline
        \textbf{Luồng chính}          &
        1. Nhân viên chọn món ăn từ danh sách thực đơn \newline
        2. Hệ thống tự động tính tổng tiền \newline
        3. Nhân viên nhập số tiền khách trả \newline
        4. Hệ thống tính tiền thừa \newline
        5. Nhân viên chọn phương thức thanh toán (Tiền mặt/VietQR) \newline
        6. Hệ thống lưu đơn hàng và cập nhật doanh thu ca                                       \\
        \hline
        \textbf{Luồng thay thế}       &
        1a. Món hết hàng: Hệ thống hiển thị thông báo và không cho phép chọn                    \\
        \hline
    \end{tabular}
\end{table}

\subsection{UC-S02: Quản lý đơn hàng}
\begin{table}[H]
    \centering
    \begin{tabular}{|p{3.5cm}|p{10cm}|}
        \hline
        \textbf{Mã UC}                & UC-S02                                                 \\
        \hline
        \textbf{Tên}                  & Quản lý đơn hàng                                       \\
        \hline
        \textbf{Tác nhân}             & Nhân viên thu ngân                                     \\
        \hline
        \textbf{Mô tả}                & Nhân viên theo dõi và cập nhật trạng thái các đơn hàng \\
        \hline
        \textbf{Điều kiện tiên quyết} & Nhân viên đã đăng nhập                                 \\
        \hline
        \textbf{Luồng chính}          &
        1. Nhân viên xem danh sách đơn hàng (Online và Tại quán) \newline
        2. Nhân viên chọn đơn hàng cần xử lý \newline
        3. Nhân viên cập nhật trạng thái (Đang chuẩn bị → Sẵn sàng → Hoàn thành) \newline
        4. Hệ thống lưu thay đổi và ghi vào lịch sử                                            \\
        \hline
    \end{tabular}
\end{table}

\subsection{UC-S03: Đóng ca làm việc}
\begin{table}[H]
    \centering
    \begin{tabular}{|p{3.5cm}|p{10cm}|}
        \hline
        \textbf{Mã UC}                & UC-S03                                         \\
        \hline
        \textbf{Tên}                  & Đóng ca làm việc                               \\
        \hline
        \textbf{Tác nhân}             & Nhân viên thu ngân                             \\
        \hline
        \textbf{Mô tả}                & Nhân viên kết thúc ca và xem báo cáo doanh thu \\
        \hline
        \textbf{Điều kiện tiên quyết} & Nhân viên có ca đang mở                        \\
        \hline
        \textbf{Luồng chính}          &
        1. Nhân viên chọn Đóng ca làm việc \newline
        2. Hệ thống hiển thị tổng kết: số đơn, doanh thu tiền mặt, doanh thu chuyển khoản \newline
        3. Nhân viên xác nhận đóng ca \newline
        4. Hệ thống lưu thông tin ca và chuyển trạng thái sang CLOSED                  \\
        \hline
    \end{tabular}
\end{table}

% ===================== QUẢN LÝ =====================
\section{Đặc tả Use Case phân hệ Quản lý}

\subsection{UC-A01: Quản lý thực đơn}
\begin{table}[H]
    \centering
    \begin{tabular}{|p{3.5cm}|p{10cm}|}
        \hline
        \textbf{Mã UC}                & UC-A01                                       \\
        \hline
        \textbf{Tên}                  & Quản lý thực đơn                             \\
        \hline
        \textbf{Tác nhân}             & Quản lý                                      \\
        \hline
        \textbf{Mô tả}                & Quản lý thêm, sửa, xóa món ăn trong thực đơn \\
        \hline
        \textbf{Điều kiện tiên quyết} & Quản lý đã đăng nhập với quyền ADMIN         \\
        \hline
        \textbf{Luồng chính}          &
        1. Quản lý xem danh sách món ăn \newline
        2. Quản lý thêm món mới (nhập tên, giá, hình ảnh, danh mục) \newline
        3. Quản lý sửa thông tin món ăn \newline
        4. Quản lý bật/tắt trạng thái Còn hàng                                       \\
        \hline
        \textbf{Luồng thay thế}       &
        2a. Thông tin không hợp lệ: Hệ thống hiển thị lỗi validation                 \\
        \hline
    \end{tabular}
\end{table}

\subsection{UC-A02: Quản lý nhân viên}
\begin{table}[H]
    \centering
    \begin{tabular}{|p{3.5cm}|p{10cm}|}
        \hline
        \textbf{Mã UC}                & UC-A02                                     \\
        \hline
        \textbf{Tên}                  & Quản lý nhân viên                          \\
        \hline
        \textbf{Tác nhân}             & Quản lý                                    \\
        \hline
        \textbf{Mô tả}                & Quản lý thêm, sửa, xóa tài khoản nhân viên \\
        \hline
        \textbf{Điều kiện tiên quyết} & Quản lý đã đăng nhập với quyền ADMIN       \\
        \hline
        \textbf{Luồng chính}          &
        1. Quản lý xem danh sách nhân viên \newline
        2. Quản lý thêm nhân viên mới (username, password, họ tên, vai trò) \newline
        3. Quản lý sửa thông tin nhân viên \newline
        4. Quản lý vô hiệu hóa tài khoản nhân viên                                 \\
        \hline
    \end{tabular}
\end{table}

\subsection{UC-A03: Quản lý nguyên liệu}
\begin{table}[H]
    \centering
    \begin{tabular}{|p{3.5cm}|p{10cm}|}
        \hline
        \textbf{Mã UC}                & UC-A03                                                \\
        \hline
        \textbf{Tên}                  & Quản lý nguyên liệu                                   \\
        \hline
        \textbf{Tác nhân}             & Quản lý                                               \\
        \hline
        \textbf{Mô tả}                & Quản lý nhập xuất kho nguyên liệu và theo dõi tồn kho \\
        \hline
        \textbf{Điều kiện tiên quyết} & Quản lý đã đăng nhập                                  \\
        \hline
        \textbf{Luồng chính}          &
        1. Quản lý xem danh sách nguyên liệu và số lượng tồn kho \newline
        2. Quản lý thêm nguyên liệu mới (tên, đơn vị, đơn giá) \newline
        3. Quản lý ghi nhận nhập kho (loại IN, số lượng, đơn giá) \newline
        4. Hệ thống tự động cập nhật số lượng tồn kho                                         \\
        \hline
        \textbf{Luồng thay thế}       &
        4a. Tồn kho dưới mức tối thiểu: Hệ thống hiển thị cảnh báo                            \\
        \hline
    \end{tabular}
\end{table}

\subsection{UC-A04: Báo cáo thống kê}
\begin{table}[H]
    \centering
    \begin{tabular}{|p{3.5cm}|p{10cm}|}
        \hline
        \textbf{Mã UC}                & UC-A04                                               \\
        \hline
        \textbf{Tên}                  & Báo cáo thống kê                                     \\
        \hline
        \textbf{Tác nhân}             & Quản lý                                              \\
        \hline
        \textbf{Mô tả}                & Quản lý xem các báo cáo doanh thu, thống kê bán hàng \\
        \hline
        \textbf{Điều kiện tiên quyết} & Quản lý đã đăng nhập                                 \\
        \hline
        \textbf{Luồng chính}          &
        1. Quản lý chọn khoảng thời gian cần xem báo cáo \newline
        2. Hệ thống hiển thị tổng doanh thu, giá vốn, lợi nhuận, số đơn hàng \newline
        3. Hệ thống hiển thị biểu đồ doanh thu theo ngày \newline
        4. Quản lý có thể xuất báo cáo ra file Excel                                         \\
        \hline
    \end{tabular}
\end{table}

\subsection{UC-A05: Quản lý đơn hàng (Admin)}
\begin{table}[H]
    \centering
    \begin{tabular}{|p{3.5cm}|p{10cm}|}
        \hline
        \textbf{Mã UC}                & UC-A05                                                 \\
        \hline
        \textbf{Tên}                  & Quản lý đơn hàng                                       \\
        \hline
        \textbf{Tác nhân}             & Quản lý                                                \\
        \hline
        \textbf{Mô tả}                & Quản lý xem và quản lý toàn bộ đơn hàng trong hệ thống \\
        \hline
        \textbf{Điều kiện tiên quyết} & Quản lý đã đăng nhập                                   \\
        \hline
        \textbf{Luồng chính}          &
        1. Quản lý xem danh sách đơn hàng với các bộ lọc \newline
        2. Quản lý lọc theo: Ngày, Trạng thái, Loại đơn, Mã đơn \newline
        3. Quản lý xem chi tiết từng đơn hàng \newline
        4. Quản lý có thể hủy đơn hàng nếu cần                                                 \\
        \hline
    \end{tabular}
\end{table}

\end{document}